\providecommand{\pdfxopt}{a-1b,mathxmp}

%\documentclass[10pt]{article}
%\usepackage{beamerarticle}
\documentclass[ignorenonframetext,slidestop]{beamer}

\input{/home/asreeve/current/class/ers102/pre_latex/luatex_pre.tex}

\setbeamertemplate{enumerate item}[square]
\setbeamertemplate{enumerate subitem}[circle]
\setbeamertemplate{enumerate subsubitem}[square]
\setbeamercolor*{enumerate item}{fg=red,bg=blue}
\setbeamercolor*{enumerate subitem}{fg=blue,bg=red}

\begin{document}

\begin{frame}{Sedimentary Rocks}
  \begin{enumerate}
  \item Sedimentary rocks form from the chemicals and particulates created during weathering of other earth materials.
  \item Sedimentary rocks form as a result of weathering, transport, and lithification.
  \item As with other rocks, the composition and texture of a sedimentary rock provides information on how and where they were formed.
  \item What you learned about minerals in the igneous rock lab is applicable to this lab.
  \end{enumerate}
\end{frame}

\begin{frame}{Sedimentary Rock Classification}
  Sedimentary rock classification is based on texture with further classification based on mineralogy and other features in the rock.
  \begin{enumerate}
  \item \textbf{Clastic} sedimentary rocks are made up of grains (solid particles).
    \begin{itemize}
    \item Shale
    \item Siltstone
    \item Sandstone
    \item Conglomerate/Breccia
    \end{itemize}
  \item \textbf{Chemical} sedimentary rocks precipitate from chemicals derived from weathering.
    \begin{itemize}
    \item Limestone (reacts with acid!)
    \item Dolostone (reacts with acid when powdered!)
    \item Rock Salt and Gypsum
    \item Chert
    \item Coal (?)
    \end{itemize}
  \end{enumerate}
\end{frame}

\begin{frame}{Sedimentary Rock and Environments}
  \begin{enumerate}
  \item  Sedimentary rocks form at or near the surface of the Earth.
  \item When considering what environment a rock formed in, use your own experience! Where have you seen the sediments in the rock?
    \begin{enumerate}
    \item Sandstone: beach, sand dune (high energy environment)
    \item Shale: deep water, flood plain (low energy environment)
    \end{enumerate}
  \end{enumerate}
\end{frame}

\begin{frame}{Doing the Lab}{Mineral ID}
  \begin{enumerate}
  \item Minerals:numbers, rocks:letters
  \item Watch the videos or examine hand samples and fill in the attached tables.
  \item More rows on the table than there are samples.
  \item Table of minerals contains minerals not in this lab, they appear in other labs
  \item Once properties defined, ID mineral sample and fill in table
  \item One new test appears in the videos, placing a drop of acid on a sample. If a sample contains calcite or similar mineral, it will fizz or effervesce.
  \end{enumerate}
The following minerals are in this lab:  calcite, quartz, halite, gypsum, kaolinite
\end{frame}

\begin{frame}{Doing the Lab}{Rock ID}
  \begin{enumerate}
  \item Group rocks by texture (clastic [and then grain size] vs. chemical)
  \item Compare with rock descriptions and ID rocks
  \end{enumerate}
The following rock samples are present: shale, sandstone, chert, conglomerate, limestone, rock salt, rock gypsum, breccia, coal
\end{frame}
\end{document}
